\section{The Single Invariant}

At a stationary path, the first variation has vanished. This does not mean
that nothing remains to be read. It means that the first readable imbalance
is gone. The next term in the action change is quadratic. It measures how
the action changes when the path is pushed in a permitted direction after
the linear debt has already been paid. This quadratic reading is the second
variation.

The second variation is not merely the next term in a formal expansion. It
is the first surviving measure of activity at stationarity. If the grammar
has already forced
\[
  \delta A[\gamma;\eta] = 0 ,
\]
then the leading nonzero response to a permitted variation has the form
\[
  \delta^2 A[\gamma;\eta,\eta].
\]
On the diagonal, where the same variation is read against itself, this is
the activity magnitude. It is the square of the energy reading in the sense
needed here: nonnegative, quadratic, and zero only when the anchored
activity itself has vanished.

The previous sections were necessary to make this sentence honest. The
stationary reader killed the first residual. The Euler--Lagrange balance
identified the local form of that killing. Spline sufficiency supplied a
finite completion that leaves no first-order knot debt. Real partitions
squeezed finite residuals toward completion. The anchor killed the nullspace
that would otherwise let nonzero activity read as zero. Only after those
obligations have been met can the second variation be called positive.

Positivity gives the invariant its force. A nonnegative quantity alone is
too weak; it may vanish on a whole quotient of unread modes. A stationary
condition alone is too weak; it may mark a saddle or a flat drift. A
completion alone is too weak; it may be smooth while carrying the wrong
activity. The anchored second variation combines the necessary features. It
is quadratic, it survives stationarity, it respects the finite completion,
and it is positive on every nonzero admissible activity.

Uniqueness follows by exclusion. A constant term cannot measure variation,
because it remains the same under every permitted slip. A linear term cannot
be the invariant at stationarity, because the grammar has forced it to
vanish. A null mode cannot be the invariant, because the anchor has pinned
it out of admissibility. A higher-order remainder cannot be the single
positive invariant, because it is not the first surviving measure and does
not supply the quadratic activity norm by which refinement is bounded. The
second variation is the only term left with the right grammar.

This is the title theorem of the chapter. The finite gauge does not merely
produce many possible quantities and then choose one by taste. It permits
finite paths, actions, variations, splines, partitions, residual bounds, and
anchors. It forbids unearned continuity, untested residue, and unpinned
null drift. Under those restrictions, the second variation is forced to be
the unique positive invariant.

The diagonal reading is especially important. A bilinear second variation
can compare two admissible variations, but the invariant appears when the
activity is compared with itself. The diagonal value is the energy square:
\[
  E(\eta)^2 = \delta^2 A[\gamma;\eta,\eta].
\]
The formula should be read grammatically. The left side names the magnitude
of admissible activity. The right side says how the finite gauge reads that
magnitude after first variation has vanished and the boundary has pinned
the nullspace. The equality identifies two renderings of one carried
quantity.

Once the invariant is available, many familiar physical readings can later
attach to it. A stiffness form can instantiate it. A Hilbert norm can be
opened from it. Charge, phase, and other labels can be tested against it.
But those later readings do not create the invariant. They inherit it. This
chapter's work is to show that the invariant is already forced by the
grammar of finite stationarity.

The single invariant is therefore not a metaphysical decoration placed on
top of measurement. It is what measurement must carry when it is finite,
gauge-bound, stationary, completed by admissible refinement, and anchored
against drift. It is the quantity that remains readable when every cheaper
reading has either vanished, been quotiented, been pinned, or been squeezed
to zero.

\section{What The Chapter Carries Forward}

The first chapter earned a place for truth to stand. This chapter has shown
what that standing place must carry when the grammar permits variation.
Finite gauge paths give measurement a way to move without abandoning its
record. Actions collect the cost of that movement. Stationarity reads the
vanishing of first residuals. Euler--Lagrange balance localizes that
vanishing. Cubic spline sufficiency gives finite anchors a balanced
completion. Real partitions squeeze supported residuals toward the completed
reading. The pinned boundary kills the nullspace. The second variation then
stands alone as the unique positive invariant.

The next chapter can therefore become concrete without changing the plot.
It will not need to invent a new invariant. It must supply a concrete
operator that carries the invariant already forced here. The finite gauge
has done its work: it has identified what must remain unchanged when
measurement is allowed to vary.
